\section*{Chapitre \ref{ch:b1:multivar} --- Fonctions de deux variables}

\begin{solution}{exo:b1:multivar:1}
$\dfrac{\partial f}{\partial x} = 2xy + y\,\eu^{xy}$,
$\dfrac{\partial f}{\partial y} = x^2 + x\,\eu^{xy}$.

$\dfrac{\partial g}{\partial x} = \dfrac{2x}{x^2+y^2}$,
$\dfrac{\partial g}{\partial y} = \dfrac{2y}{x^2+y^2}$.

$\dfrac{\partial h}{\partial x} = \dfrac{-y/x^2}{1 + y^2/x^2} =
\dfrac{-y}{x^2+y^2}$,
$\dfrac{\partial h}{\partial y} = \dfrac{x}{x^2+y^2}$.
\end{solution}

\begin{solution}{exo:b1:multivar:2}
En \omterm{prop:b1:vspaces:coordinates}{coordonnées} polaires ($\rho \to 0$)~:

$\abs{f} = \dfrac{\rho^3\abs{\cos^2\theta\sin\theta}}{\rho^2} \leq
\rho \to 0$~: \omterm{def:b1:continuity:continuous}{continue}.

$g = \cos\theta\sin\theta$~: indépendant de $\rho$, prenant des
valeurs différentes selon les demi-droites (cf.\
\cref{ex:b1:multivar:trap})~: pas de limite, non \omterm{def:b1:continuity:continuous}{continue}.

$\abs h \leq \dfrac{\rho^3(\abs{\cos^3\theta} +
\abs{\sin^3\theta})}{\rho^2} \leq 2\rho \to 0$~: \omterm{def:b1:continuity:continuous}{continue}.
\end{solution}

\begin{solution}{exo:b1:multivar:3}
Le long de $y = mx$~: $f(x, mx) = \dfrac{m x^3}{x^4 + m^2 x^2} =
\dfrac{mx}{x^2 + m^2} \to 0$ (pour $m \neq 0$~; le long de $y = 0$
et de l'axe des $y$, $f = 0$). Toute limite selon une droite vaut
donc $0$. Mais sur la parabole $y = x^2$~:
\[
f(x, x^2) = \frac{x^4}{x^4 + x^4} = \frac12 :
\]
la suite $\bigl(\frac1n, \frac{1}{n^2}\bigr) \to (0,0)$ vérifie
$f \to \frac12 \neq 0$. Non \omterm{def:b1:continuity:continuous}{continue}~: les droites ne suffisent pas
à tester les limites à deux variables.
\end{solution}

\begin{solution}{exo:b1:multivar:4}
$\frac{\partial f}{\partial x} = 3x^2y^2 + y\cos(xy)$~; puis
\[
\frac{\partial^2 f}{\partial y\,\partial x}
= 6x^2 y + \cos(xy) - xy\sin(xy) .
\]
$\frac{\partial f}{\partial y} = 2x^3 y + x\cos(xy)$~; puis
\[
\frac{\partial^2 f}{\partial x\,\partial y}
= 6x^2 y + \cos(xy) - xy\sin(xy) :
\]
égales, comme Schwarz le promet.
\end{solution}

\begin{solution}{exo:b1:multivar:5}
Par le \cref{thm:b1:multivar:chain} avec $(x, y) = (\cos t, \sin t)$~:
\[
g'(t) = -\sin t\,\frac{\partial f}{\partial x}(\cos t, \sin t)
+ \cos t\,\frac{\partial f}{\partial y}(\cos t, \sin t)
= \bigl\langle \nabla f,\ (-\sin t, \cos t)\bigr\rangle .
\]
En un extremum de $g$, $g'(t) = 0$~: $\nabla f$ est \omterm{def:b1:euclid:orthogonal}{orthogonal} au
vecteur tangent $(-\sin t, \cos t)$ du cercle, donc parallèle au
rayon vecteur $(\cos t, \sin t)$ (dans le plan, l'\omterm{def:b1:euclid:orthogonal}{orthogonal} d'un
vecteur unitaire est la droite qu'il engendre, prise
perpendiculairement). C'est l'exemple le plus simple d'un
multiplicateur de Lagrange.
\end{solution}

\begin{solution}{exo:b1:multivar:6}
$f$~: $\nabla f = (2x + y - 3,\; x + 2y) = 0$~: $y = -\frac x2$ et
$2x - \frac x2 = 3$~: $x = 2$, $y = -1$. Second ordre~: $r = 2$, $s =
1$, $t = 2$~: $rt - s^2 = 3 > 0$, $r > 0$~: minimum local (et même
global --- la fonction est quadratique) en $(2, -1)$, de valeur
$f(2,-1) = -3$.

$g$~: $\nabla g = (2x,\; -2y + 4) = 0$~: point $(0, 2)$~; $r = 2$, $s
= 0$, $t = -2$~: $rt - s^2 = -4 < 0$~: \omterm{met:b1:multivar:monge}{point col}.

$h$~: $\nabla h = \bigl(4x^3 - 4(x - y),\; 4y^3 + 4(x-y)\bigr) = 0$.
En additionnant~: $x^3 + y^3 = 0$, donc $y = -x$~; en substituant~:
$4x^3 - 8x = 0$~: $x \in \{0, \pm\sqrt2\}$. \omterm{thm:b1:multivar:critical}{Points critiques}~:
$(0,0)$, $(\sqrt2, -\sqrt2)$, $(-\sqrt2, \sqrt2)$.
En $(\pm\sqrt2, \mp\sqrt2)$~: $r = 12x^2 - 4 = 20$, $s = 4$, $t =
20$~: $rt - s^2 > 0$, $r > 0$~: minima locaux (valeur $h = 4 + 4 - 16 =
-8$). En $(0,0)$~: $r = t = -4$, $s = 4$~: $rt - s^2 = 0$~: test muet.
Examinons~: $h(x, x) = 2x^4 > 0$ et $h(x, -x) = 2x^4 - 8x^2 < 0$ pour
$x \neq 0$ petit~: les deux signes dans tout \omterm{def:b1:topology:open}{voisinage} --- pas
d'extremum à l'origine.
\end{solution}

\begin{solution}{exo:b1:multivar:7}
$x^2 - y^2$~: \omterm{def:b1:derivative:def}{dérivées} secondes $2$ et $-2$~: somme $0$. $xy$~: les
deux \omterm{def:b1:derivative:def}{dérivées} secondes pures sont nulles. $\eu^x\cos y$~:
$\frac{\partial^2}{\partial x^2} = \eu^x\cos y$,
$\frac{\partial^2}{\partial y^2} = -\eu^x\cos y$~: somme $0$.
$\ln(x^2+y^2)$~: d'après l'\cref{exo:b1:multivar:1},
\[
\frac{\partial^2}{\partial x^2}\ln(x^2+y^2)
= \frac{2(x^2+y^2) - 4x^2}{(x^2+y^2)^2}
= \frac{2(y^2 - x^2)}{(x^2+y^2)^2},
\]
et la version en $y$ en est l'opposée~: somme $0$.
\end{solution}

\begin{solution}{exo:b1:multivar:8}
\emph{\omterm{def:b1:linmaps:projection}{Projection}~:} le plan $P$ a pour normale $n = (1, 2, -1)$ et
passe par $A = (4, 0, 0)$. Le point le plus proche de l'origine est
le projeté de $O$~: $p = O + \frac{\langle A - O, n\rangle}{\norm
n^2}\,n = \frac{4}{6}(1,2,-1) = \bigl(\frac23, \frac43,
-\frac23\bigr)$, à distance $\frac{4}{\sqrt 6}$. (Formule~: la
distance de l'origine au plan $\langle X, n \rangle = c$ vaut
$\frac{\abs c}{\norm n}$ avec $c = 4$.)

\emph{Minimisation~:} $f(x,y) = x^2 + y^2 + (x + 2y - 4)^2$. En
posant $w = x + 2y - 4$ pour le dernier facteur~:
\[
\nabla f = \bigl(2x + 2w,\; 2y + 4w\bigr) = 0
\iff x = -w \text{ et } y = -2w .
\]
En reportant dans la définition de $w$~: $w = -w - 4w - 4$, donc $w =
-\frac23$, d'où $x = \frac23$, $y = \frac43$, et $z = w =
-\frac23$. Même point que par la \omterm{def:b1:linmaps:projection}{projection}, à distance
$\sqrt{\frac49 + \frac{16}{9} + \frac49} = \frac{2\sqrt 6}{3} =
\frac{4}{\sqrt6}$~; c'est un minimum puisque $f \to \infty$ à
l'infini (une quadratique définie positive plus des termes
\omterm{def:b1:linmaps:def}{linéaires}).
\end{solution}

\begin{solution}{exo:b1:multivar:9}
Avec $\rho = \sqrt{x^2+y^2}$~: $\frac{\partial \rho}{\partial x} =
\frac{x}{\rho}$, donc
\[
\frac{\partial f}{\partial x} = \varphi'(\rho)\,\frac{x}{\rho},
\qquad
\frac{\partial^2 f}{\partial x^2}
= \varphi''(\rho)\,\frac{x^2}{\rho^2}
+ \varphi'(\rho)\,\Bigl(\frac{1}{\rho} -
\frac{x^2}{\rho^3}\Bigr),
\]
(règles de dérivation du quotient et des fonctions composées). En
ajoutant l'expression symétrique en $y$~: les termes en $\varphi''$
donnent $\frac{x^2 + y^2}{\rho^2} = 1$, les termes en $\varphi'$
donnent $\frac{2}{\rho} - \frac{x^2 +
y^2}{\rho^3} = \frac{1}{\rho}$~:
\[
\Delta f = \varphi'' + \frac{\varphi'}{\rho} .
\]
Harmoniques radiales~: résolvons $\varphi'' + \frac{\varphi'}{\rho}
= 0$~: l'équation du premier ordre $u' + \frac u\rho = 0$ pour $u =
\varphi'$ donne $u = \frac{c}{\rho}$
(\cref{thm:b1:diffeq:homogeneous1}), puis $\varphi = c\ln\rho + d$.
Les fonctions harmoniques radiales sont les
$c\,\ln\sqrt{x^2 + y^2} + d$ --- le potentiel logarithmique de
l'\cref{exo:b1:multivar:7} et les constantes, rien d'autre.
\end{solution}

\begin{solution}{exo:b1:multivar:10}
Pour $z = xy$ en $(1,1)$~: \omterm{def:b1:multivar:partial}{dérivées partielles} $y = 1$ et $x = 1$,
plan tangent $z = 1 + (x - 1) + (y - 1) = x + y - 1$. Un plan
tangent horizontal signifie que les deux \omterm{def:b1:multivar:partial}{dérivées partielles}
s'annulent, c'est-à-dire $\nabla f = 0$~: d'après
l'\cref{ex:b1:multivar:extrema}, exactement les \omterm{thm:b1:multivar:critical}{points critiques}
$(0, 0)$ et $(1, 1)$, avec les plans horizontaux $z = 0$ et $z =
-1$. «~Plan tangent horizontal~» et «~\omterm{thm:b1:multivar:critical}{point critique}~» sont la
même notion, vue sur le graphe et dans la formule.
\end{solution}

\begin{solution}{exo:b1:multivar:11}
\begin{enumerate}
  \item À $y$ fixé, la fonction d'une variable $x \mapsto
        f(x,y)$ a une \omterm{def:b1:derivative:def}{dérivée} nulle sur $\R$, donc est constante
        (\cref{thm:b1:derivative:mvt})~: $f(x, y) = f(0, y)$ pour
        tout $x$~: $f$ ne dépend que de $y$.
  \item Posons $g(u, v) = f\bigl(\frac{u+v}2, \frac{u-v}2\bigr)$.
        Par la \omterm{thm:b1:derivative:operations}{règle de la chaîne}
        (\cref{thm:b1:multivar:chain}, appliquée en la variable
        $v$ à $u$ figé)~:
        \[
        \frac{\partial g}{\partial v}
        = \frac12\,\frac{\partial f}{\partial x}
        - \frac12\,\frac{\partial f}{\partial y} = 0 .
        \]
        D'après (1), $g$ ne dépend que de $u$~: $g(u, v) =
        \varphi(u)$ avec $\varphi$ de classe $C^1$, et en
        inversant le \omterm{thm:b1:integration:parts}{changement de variables} ($u = x + y$, $v = x
        - y$),
        \[
        f(x, y) = \varphi(x + y).
        \]
        Réciproquement, toute telle $f$ vérifie $f_x = f_y =
        \varphi'$~: les solutions sont exactement les fonctions de
        classe $C^1$ de $x + y$.
\end{enumerate}
\end{solution}

\begin{solution}{exo:b1:multivar:12}
Sur le disque \omterm{def:b1:topology:open}{ouvert}, un extremum serait \omterm{thm:b1:multivar:critical}{critique}~: $\nabla f = (y,
x) = 0$ seulement en $(0,0)$, où $f = 0$~; c'est un \omterm{met:b1:multivar:monge}{point col} ($f(\pm
\varepsilon, \pm\varepsilon) = \varepsilon^2 > 0 > -\varepsilon^2
= f(\pm\varepsilon, \mp\varepsilon)$), donc pas d'extremum là. Par
le théorème admis des bornes atteintes, les bornes sont atteintes,
nécessairement sur le cercle \omterm{def:b1:topology:closure}{frontière}. Là, avec
l'\cref{exo:b1:multivar:5},
\[
f(\cos t, \sin t) = \cos t\sin t = \frac{\sin 2t}{2}
\in \intcc{-\tfrac12}{\tfrac12},
\]
avec maximum $\frac12$ en $t = \frac\pi4, \frac{5\pi}4$ (points
$\pm\frac{1}{\sqrt2}(1,1)$) et minimum $-\frac12$ en $t =
\frac{3\pi}4, \frac{7\pi}4$ (points $\pm\frac{1}{\sqrt2}(1,-1)$).
Maximum global $\frac12$, minimum global $-\frac12$.
\end{solution}

\begin{solution}{pb:b1:multivar:1}
\textbf{1.} En développant le membre de droite~:
\[
r\Bigl(h + \frac srk\Bigr)^{\!2} + \frac{rt - s^2}{r}k^2
= rh^2 + 2shk + \frac{s^2}{r}k^2 + \frac{rt - s^2}{r}k^2
= q(h,k) .
\]
Si $rt - s^2 > 0$~: les deux carrés portent le signe du facteur
$r$, et $q(h,k) = 0$ force $k = 0$ puis $h = 0$~: signe strict de
$r$ hors de l'origine. Si $rt - s^2 < 0$~: $q(1, 0) = r$ tandis que
$q\bigl(-\frac sr, 1\bigr) = \frac{rt - s^2}{r}$ est de signe
opposé~: les deux signes apparaissent.

\textbf{2.} Si $r = 0 \neq t$, on échange les rôles de $h$ et $k$
(la forme canonique en $k$), avec les mêmes conclusions~;
remarquons que $rt - s^2 = -s^2 < 0$ dès que $s \neq 0$, et de fait
$q = 2shk + tk^2$ prend alors les deux signes ($k$ petit devant
$h$). Si $r = t = 0$~: $q = 2shk$, deux signes si et seulement si
$s \neq 0$, c'est-à-dire si et seulement si $rt - s^2 = -s^2 < 0$.
Bilan~: $q$ prend les deux signes $\iff rt - s^2 < 0$~; $q$ ne
s'annule qu'à l'origine (forme définie) $\iff rt - s^2 > 0$ (ce
qui force $r \neq 0$, puisque $r = 0$ donne $q(1,0) = 0$).

\textbf{3.} Avec $f_x = rx + sy + \beta$ et $f_y = sx + ty +
\gamma$, le développement direct donne
\[
f(a + h, b + k) = f(a, b) + h\,f_x(a,b) + k\,f_y(a,b)
+ \tfrac12 q(h, k),
\]
sans termes d'ordre supérieur (la fonction est un \omterm{def:b1:poly:def}{polynôme} de
degré $2$). En un \omterm{thm:b1:multivar:critical}{point critique} la partie \omterm{def:b1:linmaps:def}{linéaire} s'annule~:
$f(X_0 + H) - f(X_0) = \frac12 q(H)$ \emph{exactement}, de sorte
que l'étude de signe des questions 1 et 2 classe~: minimum global
strict ($rt - s^2 > 0$, $r > 0$), maximum global strict ($rt - s^2
> 0$, $r < 0$), \omterm{met:b1:multivar:monge}{point col} --- les deux signes dans tout \omterm{def:b1:topology:open}{voisinage}
--- lorsque $rt - s^2 < 0$. Cela démontre
la \cref{met:b1:multivar:monge} pour les fonctions quadratiques.

\textbf{4.} La forme $q - c(h^2 + k^2)$ a pour données $r - c$, $s$,
$t - c$. Avec $c = \frac{rt - s^2}{r + t}$ (noter que $t > 0$
puisque $rt > s^2 \geq 0$ et $r > 0$)~:
\[
(r - c)(t - c) - s^2 = rt - s^2 - c(r + t) + c^2 = c^2 \geq 0,
\]
et $r - c \geq 0$ ($c \leq r \iff rt - s^2 \leq r^2 + rt$,
vrai). Si $r - c > 0$, la forme canonique de la question 1 montre
que $q - c(h^2 + k^2) \geq 0$~; si $r - c = 0$, alors $s = 0$ (car
$-s^2 = -c^2 \leq \dots$ force la quantité affichée à être
$\geq 0$ avec un premier facteur nul) et la forme vaut $(t - c)k^2
\geq 0$. Dans les deux cas $q(h,k) \geq c(h^2 + k^2)$.

\textbf{5.} D'après les questions 3 et 4, $f(X) = f(X_0) + \frac12
q(X - X_0) \geq f(X_0) + \frac c2\norm{X - X_0}^2 \to +\infty$
quand $\norm X \to \infty$~: $f$ est coercive, et l'inégalité est
stricte pour $X \neq X_0$~: le \omterm{thm:b1:multivar:critical}{point critique} est l'unique
minimiseur global. Pour la cubique $f = x^3 + y^3 - 3xy$ de
l'\cref{ex:b1:multivar:extrema}, aucune telle conclusion ne vaut~:
$f(x,0) = x^3 \to -\infty$, et le minimum local en $(1,1)$ n'est
pas global --- c'est l'exactitude du développement quadratique qui
faisait défaut.

\textbf{6.} En développant le carré de la norme~:
\[
f(u,v) = u^2\norm{C_1}^2 + 2uv\,\langle C_1, C_2\rangle +
v^2\norm{C_2}^2 - 2u\,\langle C_1, b\rangle - 2v\,\langle C_2,
b\rangle + \norm b^2 ,
\]
fonction quadratique de $(u, v)$ dont la partie quadratique est
$\frac12(ru^2 + 2suv + tv^2)$ avec $r = 2\norm{C_1}^2$, $s =
2\langle C_1, C_2\rangle$, $t = 2\norm{C_2}^2$.

\textbf{7.} $rt - s^2 = 4\bigl(\norm{C_1}^2\norm{C_2}^2 -
\langle C_1, C_2\rangle^2\bigr) \geq 0$ par Cauchy--Schwarz
(\cref{thm:b1:euclid:cs}), avec égalité exactement quand $C_1, C_2$
sont proportionnels (ou l'un des deux nul), c'est-à-dire quand la
famille est liée. Donc $rt - s^2 > 0 \iff (C_1, C_2)$ \omterm{def:b1:vspaces:free}{libre}.

\textbf{8.} $\nabla f = 0$ s'écrit
\[
\norm{C_1}^2 u + \langle C_1, C_2\rangle v = \langle C_1,
b\rangle,
\qquad
\langle C_1, C_2\rangle u + \norm{C_2}^2 v = \langle C_2,
b\rangle,
\]
les \omterm{pb:b1:multivar:1}{équations normales}, avec la matrice de Gram à gauche~; elles
disent $\langle uC_1 + vC_2 - b,\ C_i\rangle = 0$ pour $i = 1, 2$,
c'est-à-dire $b - p \perp \operatorname{Vect}(C_1, C_2)$ pour $p =
uC_1 + vC_2$. Par la partie I (questions 3 et 5) l'unique \omterm{thm:b1:multivar:critical}{point
critique} est l'unique minimiseur global, et par
le \cref{thm:b1:euclid:projection} la caractérisation «~$p \in
F$, $b - p \perp F$~» identifie $p$ à la \omterm{thm:b1:euclid:projection}{projection orthogonale}
de $b$ sur $F = \operatorname{Vect}(C_1, C_2)$.

\textbf{9.} $\norm{b - p}^2 = \norm b^2 - 2\langle b, p\rangle +
\norm p^2$, et $\langle p, b - p\rangle = 0$ donne $\norm p^2 =
\langle p, b\rangle$~: le minimum vaut $\norm b^2 - \langle p,
b\rangle$. Pythagore~: $\norm b^2 = \norm p^2 + \norm{b - p}^2$
--- le vecteur de données se scinde en sa part expliquée $p$ et sa
part résiduelle $b - p$, \omterm{def:b1:euclid:orthogonal}{orthogonales} l'une à l'autre.

\textbf{10.} $E(\alpha, \beta) = \norm{\alpha C_1 + \beta C_2 -
b}^2$ avec les $C_1, C_2, b$ indiqués~; les \omterm{pb:b1:multivar:1}{équations normales} de
la question 8 s'écrivent, terme à terme,
\[
n\alpha + \Bigl(\sum x_i\Bigr)\beta = \sum y_i,
\qquad
\Bigl(\sum x_i\Bigr)\alpha + \Bigl(\sum x_i^2\Bigr)\beta = \sum
x_i y_i .
\]

\textbf{11.} En divisant par $n$~: $\alpha + \beta\bar x = \bar y$
et $\alpha\bar x + \beta\bigl(v_x + \bar x^2\bigr) = c_{xy} +
\bar x\bar y$. En substituant $\alpha = \bar y - \beta\bar x$ dans
la seconde~: $\beta v_x = c_{xy}$, donc
\[
\beta = \frac{c_{xy}}{v_x}, \qquad \alpha = \bar y - \beta\bar
x ,
\]
et la première équation dit précisément que $(\bar x, \bar y)$
est sur la droite.

\textbf{12.} $v_x = \frac1n\sum(x_i - \bar x)^2 \geq 0$, nul si et
seulement si tous les $x_i$ valent $\bar x$. Et $rt - s^2 =
4\bigl(n\sum x_i^2 - (\sum x_i)^2\bigr) = 4n^2 v_x$~: la condition
de liberté de la question 7 est $v_x > 0$, c'est-à-dire les $x_i$
non tous égaux.

\textbf{13.} Avec $\alpha = \bar y - \beta\bar x$, centrons les
données ($\tilde x_i = x_i - \bar x$, $\tilde y_i = y_i - \bar
y$)~:
\[
E(\alpha, \beta) = \sum_i(\tilde y_i - \beta\tilde x_i)^2
= n\,v_y - 2\beta\,n\,c_{xy} + \beta^2 n\,v_x ,
\]
minimisée en $\beta = c_{xy}/v_x$, de valeur $E_{\min} =
n\bigl(v_y - \frac{c_{xy}^2}{v_x}\bigr) = n v_y(1 - \rho^2)$.
Comme $E_{\min} \geq 0$ et $v_y > 0$~: $\rho^2 \leq 1$, et
$\abs\rho = 1$ si et seulement si $E_{\min} = 0$, c'est-à-dire si
et seulement si tous les résidus sont nuls~: les points sont
exactement sur la droite.

\textbf{14.} $n = 4$~: $\bar x = \frac64 = 1.5$, $\bar y = 2$,
$\sum x_i^2 = 14$ donc $v_x = 3.5 - 2.25 = 1.25$, $\sum x_iy_i =
0 + 1 + 6 + 12 = 19$ donc $c_{xy} = 4.75 - 3 = 1.75$. Ainsi
\[
\beta = \frac{1.75}{1.25} = 1.4,
\qquad
\alpha = 2 - 1.4\times1.5 = -0.1 :
\qquad y = 1.4\,x - 0.1 .
\]
Valeurs ajustées $-0.1,\ 1.3,\ 2.7,\ 4.1$~; résidus $0.1,\ -0.3,\
0.3,\ -0.1$~; $E_{\min} = 0.01 + 0.09 + 0.09 + 0.01 = 0.2$
(vérification~: $v_y = \frac{26}4 - 4 = 2.5$ et $4(2.5 -
\frac{1.75^2}{1.25}) = 4(2.5 - 2.45) = 0.2$).

\textbf{15.} Les deux \omterm{pb:b1:multivar:1}{équations normales} de la question 10 sont
exactement $\sum_i\varepsilon_i = 0$ et $\sum_i x_i\varepsilon_i =
0$~: le vecteur des résidus est \omterm{def:b1:euclid:orthogonal}{orthogonal} à $C_1 = (1, \dots, 1)$
et à $C_2 = (x_i)$ --- à tout l'espace du modèle. En particulier
la droite optimale équilibre toujours ses erreurs~: leur somme est
nulle.

\textbf{16.} $\frac{\dd}{\dd\alpha}\sum(y_i - \alpha)^2 =
-2\sum(y_i - \alpha) = 0$ donne $\alpha = \bar y$ (et la \omterm{def:b1:derivative:def}{dérivée}
seconde $2n > 0$ en fait le minimum global, la fonction étant une
quadratique coercive d'une variable). Valeur minimale
$\sum(y_i - \bar y)^2 = n v_y$~: la variance est l'erreur
quadratique incompressible d'un modèle constant.

\textbf{17.} $\frac{\dd}{\dd\beta}\sum(y_i - \beta x_i)^2 =
-2\sum x_i(y_i - \beta x_i) = 0$ donne $\beta_0 = \frac{\sum
x_iy_i}{\sum x_i^2}$. En écrivant les deux pentes sur des
quantités centrées~: $\beta_0 = \frac{c_{xy} + \bar x\bar y}{v_x +
\bar x^2}$ vaut $\frac{c_{xy}}{v_x}$ si et seulement si
$c_{xy}\bar x^2 = \bar x\bar y\,v_x$, c'est-à-dire si et seulement
si $\bar x = 0$ (abscisses centrées) ou $\bar y = \beta\bar x$ ---
ce dernier cas signifiant $\alpha = 0$~: la \omterm{pb:b1:multivar:1}{droite de régression}
complète passe déjà par l'origine.

\textbf{18.} $g(\tau) = \norm{M - P_0}^2 - 2\tau\langle M - P_0,
w\rangle + \tau^2$ est minimale en $\tau^* = \langle M - P_0,
w\rangle$, de valeur $d(M,D)^2 = \norm{M - P_0}^2 - \langle M -
P_0, w\rangle^2$. Dans le plan, complétons $w$ en une \omterm{def:b1:vspaces:free}{base}
\omterm{def:b1:euclid:orthogonal}{orthonormale} $(w, n)$ avec $n = \frac{(a, b)}{\sqrt{a^2+b^2}}$
(normale unitaire de $D$)~: alors $\norm{M - P_0}^2 = \langle M -
P_0, w\rangle^2 + \langle M - P_0, n\rangle^2$, donc $d(M, D) =
\abs{\langle M - P_0, n\rangle}$, et avec $aP_{0x} + bP_{0y} =
-c$~:
\[
d = \frac{\abs{a x_0 + b y_0 + c}}{\sqrt{a^2 + b^2}} .
\]

\textbf{19.} De $E_{\min}/n = v_y - \frac{c_{xy}^2}{v_x}$ et
$\beta^2 v_x = \frac{c_{xy}^2}{v_x}$~:
\[
v_y = \beta^2 v_x + \frac{E_{\min}}n :
\]
variance totale $=$ variance le long de la droite ajustée $+$
variance résiduelle. (En divisant par $v_y$~: $1 = \rho^2 + (1 -
\rho^2)$, la part de variance «~expliquée~» par la droite est
$\rho^2$.)

\textbf{20.} L'espace du modèle est
$\operatorname{Vect}\bigl(C_1, C_2, C_3\bigr)$ avec $C_1 = (1)$,
$C_2 = (x_i)$, $C_3 = (x_i^2)$~; minimiser $\norm{aC_1 + bC_2 +
cC_3 - b}^2$ conduit, exactement comme à la question 8, au système
de Gram $3\times3$, dont la matrice a pour coefficients $\langle
C_i, C_j\rangle = \sum_k x_k^{\,i+j-2}$~: la matrice des moments
affichée. Elle est inversible si et seulement si $(C_1, C_2, C_3)$
est \omterm{def:b1:vspaces:free}{libre} (\cref{exo:b1:euclid:11})~; une relation $a + bx_k +
cx_k^2 = 0$ pour tout $k$ fait de chaque $x_k$ une racine d'un même
\omterm{def:b1:poly:def}{polynôme} de degré $\leq 2$, impossible avec trois valeurs
distinctes à moins que $a = b = c = 0$. Comparer avec les matrices
des moments du devoir maison du \cref{ch:b1:det}, dont les
déterminants étaient des carrés de valeurs de Vandermonde.

\textbf{21.} Première~: $r = 2, s = 1, t = 2$, $rt - s^2 = 3 > 0$,
$r > 0$~: minimum global strict à l'origine. Deuxième~: $r = 2, s
= 3, t = 2$, $rt - s^2 = -5 < 0$~: \omterm{met:b1:multivar:monge}{point col}. Troisième~: $r = s = t =
2$, $rt - s^2 = 0$~: dégénérée --- mais $x^2 + 2xy + y^2 = (x +
y)^2 \geq 0$ s'annule sur toute la droite $y = -x$~: un minimum
global (non strict) atteint le long d'une droite, invisible au
test du déterminant.

\textbf{22.} Nouvelles sommes ($n = 5$)~: $\bar x = 3.2$, $\bar y =
1.6$, $\sum x_iy_i = 19$ donc $c_{xy} = 3.8 - 5.12 = -1.32$,
$\sum x_i^2 = 114$ donc $v_x = 22.8 - 10.24 = 12.56$. Nouvelle
pente~:
\[
\beta = \frac{-1.32}{12.56} \approx -0.105 :
\]
un seul point a transformé une tendance nettement croissante
($\beta = 1.4$) en une tendance légèrement décroissante. L'erreur
quadratique facture un résidu $\varepsilon^2$, de sorte qu'un seul
point lointain --- $\abs{x_{i} - \bar x}$ grand \emph{et} résidu
grand --- domine à la fois $c_{xy}$ et $v_x$~: les \omterm{pb:b1:multivar:1}{moindres carrés}
sont efficaces, mais pas robustes.

\textbf{23.} Posons $W = \sum w_i$ et définissons $\bar x_w =
\frac1W\sum w_ix_i$, $\bar y_w$ de même, $v_x^w = \frac1W\sum
w_ix_i^2 - \bar x_w^2$, $c^w_{xy} = \frac1W\sum w_ix_iy_i - \bar
x_w\bar y_w$. L'\omterm{def:b1:logic:map}{application} $\langle u, v\rangle_w = \sum_i
w_iu_iv_i$ est un \omterm{def:b1:euclid:def}{produit scalaire} sur $\R^n$ ($w_i > 0$ donne le
caractère défini), donc la partie II s'applique mot pour mot et
les \omterm{pb:b1:multivar:1}{équations normales}, divisées par $W$, deviennent
\[
\alpha + \beta\bar x_w = \bar y_w,
\qquad
\alpha\bar x_w + \beta(v^w_x + \bar x_w^2) = c^w_{xy} + \bar
x_w\bar y_w ,
\]
d'où $\beta = c^w_{xy}/v^w_x$ et $\alpha = \bar y_w -
\beta\bar x_w$~: les mêmes formules, avec toutes les moyennes
pondérées.

\textbf{24.} $E_{\min} = \sum\varepsilon_i^2 = 0$ force chaque
$\varepsilon_i = 0$~: tous les points sont exactement sur la
droite~; par la question 13 c'est le cas $\abs\rho = 1$. Test~:
pour $(1,1), (2,3), (3,5)$~: $\bar x = 2$, $\bar y = 3$, $v_x =
\frac{14}3 - 4 = \frac23$, $c_{xy} = \frac{22}3 - 6 = \frac43$~:
$\beta = 2$, $\alpha = -1$, et de fait $y_i = 2x_i - 1$ pour les
trois points~; $v_y = \frac{35}3 - 9 = \frac83$ et $\rho^2 =
\frac{(4/3)^2}{(2/3)(8/3)} = 1$.

\textbf{25.} (i) Pour les fonctions quadratiques le développement
à l'ordre deux est une \emph{identité}, de sorte que l'étude de
signe de la forme quadratique --- de l'algèbre pure, questions 1
et 2 --- classe les \omterm{thm:b1:multivar:critical}{points critiques} sans aucun théorème de Taylor
admis. (ii) Les \omterm{pb:b1:multivar:1}{équations normales} disent «~résidu \omterm{def:b1:euclid:orthogonal}{orthogonal} à
l'espace du modèle~», donc le minimum analytique \emph{est} la
\omterm{thm:b1:euclid:projection}{projection orthogonale} du vecteur de données~: l'analyse et la
géométrie euclidienne calculent le même objet. (iii) La
corrélation $\rho = c_{xy}/\sqrt{v_xv_y}$ mesure la part $\rho^2$
de la variance expliquée par la droite, et Cauchy--Schwarz la
borne~: $\abs\rho \leq 1$, avec égalité seulement pour des données
alignées. (iv) \omterm{def:b1:vspaces:free}{Bases} et liberté (chapitre 18), matrices de Gram et
des moments et leurs déterminants (chapitres 22 et 23), et
\omterm{thm:b1:euclid:projection}{projection orthogonale} (chapitre 23) ont tous reparu comme les
pièces d'un même algorithme. Les parties II et III développent la
\emph{méthode des \omterm{pb:b1:multivar:1}{moindres carrés}} (Legendre, Gauss) à travers ses
\emph{\omterm{pb:b1:multivar:1}{équations normales}}.
\end{solution}
